\section{Related Work}
\label{sec:related}

\noindent\textbf{Object hallucination detection in LVLMs.}
Object hallucination was formalized as an image-captioning failure by CHAIR \cite{rohrbach2018chair} and has since become a central reliability problem for LVLMs. Recent detectors differ in what internal signal they treat as evidence. Uncertainty methods use the output distribution, such as token likelihood or entropy. Representation methods, such as IC, use internal visual representations to estimate whether the object is supported \cite{jiang2024interpreting}. Attention-based methods instead score how object tokens route computation: SVAR uses object-to-image attention mass \cite{jiang2024devils}, PAS measures object-token attention to preliminary generated tokens \cite{hoangxuan2026pas}, and recent fine-grained grounding scores combine patch-level dispersion and consistency \cite{beyond2026global}. We use these detectors as baselines for a stricter selectivity question, then test TDEV as a target-evidence alternative: does the signal remain selective for object evidence after controlling for generation position and object identity?

\noindent\textbf{Training-free hallucination mitigation.}
Training-free mitigation methods intervene at several different points. Visual Contrastive Decoding contrasts logits from original and visually distorted inputs to reduce language-prior reliance \cite{leng2023vcd}, while OPERA penalizes over-trust in summary tokens and reallocates decoding when that pattern appears \cite{huang2023opera}. Attention-side methods act more directly on routing: PAI increases visual attention and combines the attention branch with contrastive/CFG-style decoding in the full method \cite{liu2024pai}, ClearSight amplifies mid-layer image signals and suppresses prefix influence \cite{yin2025clearsight}, Visual Attention Sink redistributes mass away from high-activation sink tokens toward visual tokens \cite{kang2025visualsink}, and DAMRO suppresses high-attention outlier/background visual tokens \cite{gong2024damro}. Recent caption-query and head-steering methods, including CAI and CAST, use caption-induced attention patterns or steering vectors to strengthen visual perception during inference \cite{li2025cai,li2026cast}; related activation, phase-aware, or region-aware steering methods similarly recalibrate hidden states, visual tokens, heads, or regional attention budgets \cite{yin2026dynamic,kim2026focus,xu2026regionaware}. These works are closest to the practical mitigation direction, but they still motivate the same evaluation question we study: does stronger visual routing verify the queried target under semantic-neighbor negatives? Post-hoc revision methods such as LURE, Woodpecker, and Volcano correct generated descriptions using co-occurrence/uncertainty/position factors, visual validation, or self-feedback \cite{zhou2023lure,yin2023woodpecker,lee2023volcano}. Our experiments focus on attention-intervention components and TDEV gates because they directly test the allocation--verification gap; decoding, steering, and post-hoc methods are complementary baselines for broader end-to-end comparisons.

\noindent\textbf{Evaluation shortcuts and evidence verification.}
Several recent analyses show that LVLM hallucination benchmarks can reward shortcut behavior. The Mirage study argues that contrastive-decoding gains on POPE can come from output-distribution shifts rather than genuine mitigation \cite{yin2025mirage}. AdaVIB connects hallucination to overconfidence in irrelevant visual features after visual tokens enter the language embedding space \cite{bai2025adavib}. These findings motivate a stricter question than whether the model attends to the image: does the attended evidence verify the target object claim? Our diagnostic protocol extends this question to both detection and mitigation by combining position control, answer-prior control, attention-shift selectivity, and semantic-neighbor negatives.

\noindent\textbf{Open-vocabulary region evidence.}
Open-vocabulary detectors such as OWLv2 localize text-specified object categories without retraining for the closed set used by a downstream hallucination benchmark \cite{minderer2023scaling}. We use OWLv2 as an external positive control rather than as a replacement for LVLM grounding: if target-discriminative region evidence passes our controls while decoder-attention mass fails them, the gap is evidence about the proxy, not evidence that visual grounding is impossible.
